Cryptocurrency arbitrage has been studied from multiple perspectives, including graph-theoretic opportunity detection, system-level arbitrage platforms, and execution strategy optimization. Most prior work focuses primarily on identifying arbitrage opportunities, while comparatively less attention has been devoted to modeling execution feasibility under real-world constraints. This section situates our approach within these research streams.

\subsection{Graph-Theoretic Arbitrage Detection (Negative Cycles, Bellman--Ford)}

The dominant formulation of arbitrage detection models financial markets as directed weighted graphs. Nodes represent currencies or tokens, and exchange rates are encoded as edges. Applying a logarithmic transformation,
\[
w(e) = -\log(r(e)),
\] %this one can be inline
converts multiplicative exchange rates into additive costs, allowing arbitrage opportunities to be identified as negative-weight cycles. This transformation enables the use of classical shortest-path algorithms such as Bellman--Ford and its variants.

Oantă and Coroiu (2023) adopt this paradigm for cross-exchange cryptocurrency markets in their system \emph{Crypto Advisor}, which aggregates real-time exchange data and detects arbitrage cycles using graph-based methods. The primary contribution is infrastructural, providing a scalable web-based platform for reporting opportunities. However, the detection framework assumes idealized execution conditions, including sufficient liquidity, static prices during execution, and negligible operational risk.

Similar graph-based formulations appear in decentralized finance (DeFi) arbitrage systems, improved negative-cycle detection algorithms, and real-time high-frequency frameworks. Across these works, the core objective is opportunity existence detection rather than execution-aware feasibility analysis.

\subsection{Pathfinding Approaches in Financial Routing}

Beyond negative-cycle detection, shortest-path formulations have been proposed to compute cost-minimizing trading routes. Algorithms such as Dijkstra’s method and its weighted variants provide efficient routing in additive cost graphs. However, these approaches typically assume static and deterministic edge weights.

In cross-exchange contexts, routing decisions involve heterogeneous constraints, including withdrawal fees, transfer costs, and asset availability across exchanges. Traditional shortest-path algorithms can model additive costs but do not incorporate domain-specific risk factors or execution uncertainty during search. As a result, routing decisions often remain detached from real-world feasibility.

\subsection{Heuristic Search in Constrained Optimization}

Heuristic search methods such as A* have been widely studied in artificial intelligence for solving constrained pathfinding and planning problems. By incorporating domain-specific heuristics, A* can guide search toward promising regions of the state space while preserving optimality under admissible conditions.

Although heuristic search has been extensively applied in robotics, logistics, and multi-agent path finding (MAPF), its application to cross-exchange cryptocurrency arbitrage remains limited. Prior financial graph formulations focus on detecting negative cycles or computing shortest paths without embedding execution risk directly into heuristic guidance. This creates an opportunity to apply AI planning techniques to arbitrage execution planning under real-world constraints.

\subsection{Arbitrage Under Latency, Fees, and Market Frictions}

Execution feasibility has been studied more extensively in single-exchange optimal execution literature, often using reinforcement learning or stochastic control methods (e.g., Kearns \& Nevmyvaka, 2006). These approaches optimize trade scheduling within a single market, accounting for price impact and liquidity dynamics.

However, cross-exchange arbitrage introduces additional constraints, including blockchain transfer latency, withdrawal fees, liquidity fragmentation, and exchange reliability. Stablecoin markets, in particular, exhibit narrow spreads and frequent micro-inefficiencies, making them highly sensitive to execution frictions. Existing detection-oriented methods do not explicitly integrate these constraints into the search process, increasing the risk of theoretical but non-executable opportunities.

\subsection{Comparison to System-Oriented Arbitrage Platforms}

Several systems-oriented contributions focus on scalable architectures for aggregating exchange data, computing spreads, and streaming arbitrage signals through APIs or dashboards. These works emphasize infrastructure, latency reduction, and data engineering robustness.

While valuable from a systems perspective, such platforms typically treat arbitrage computation as a spread comparison or post-hoc cycle detection task. Execution constraints—including liquidity depth, slippage, transfer delays, and operational risk—are evaluated after opportunity detection rather than embedded within the search itself.

In contrast, our work integrates feasibility modeling directly into the search process. Instead of first detecting theoretical arbitrage and then evaluating execution risk, we formulate arbitrage planning as a constraint-aware search problem and apply A* with domain-specific heuristics. This approach bridges opportunity detection and executable route construction, particularly suited to the low-margin, friction-sensitive context of stablecoin markets.
